\section{Discussion}

The experiments support a two-stage design test for context-routing systems.
The first stage measures structural headroom: whether state conditioning
creates value that a perfect standalone ranking cannot reach. The second stage
measures how accurately marginal utility can be estimated from information
available to the router. Tasks with little headroom need no sequential
marginal routing. Tasks with substantial headroom require predictor-response
information and enough decision margin for the router to distinguish
candidates.

The multi-target study shows why both stages matter. Structural headroom is
positive on every evaluated target, but learned routing is more reliable on
some targets than others. METR-LA has the widest target-level variation. This
pattern matches the mechanism analysis: structural opportunity alone does not
determine how much value a router can realize. Predictor-response features and
the utility gap between the best candidates are the operating quantities.

The main limitations are computational and structural. R-MUR is greedy and can
miss pure complementarity, where every individual candidate has nonpositive
utility but a bundle is useful. Its labels require counterfactual evaluations
of the fixed predictor, so offline label cost is part of the system cost. The
calibration-derived stopping threshold controls the gain--harm trade-off, and
its guarantee is per-decision.

\subsection{Structural headroom is not the same as realizable headroom}

The strong-backbone study separates two quantities that the earlier results
could not distinguish. Structural headroom is a property of the predictor and
the candidate pool: with the subset-capable transformer, the sequential oracle
still beats the standalone oracle by $.0148$ absolute on METR-LA, so state
conditioning continues to create value. Realizable headroom is a property of
the information available to the router. With the ridge expert the learned
router converted a substantial share of the structural opportunity into gain;
with the nonlinear expert the response-free routers are statistically
indistinguishable from random selection even though the same structural
opportunity is present.

The general-loss analysis of \cref{sec:analysis} explains the mechanism. The
expectation of the exact identity conditions on the state as
$\E[m(j\mid A)\mid x,A]=\langle g_A(x),d_{A,j}\rangle-\|d_{A,j}\|^2-\lambda
c_j$ with $g_A(x)=2(\mu(x)-f_A(x))$. The second term is known as soon as the
predictor has been probed, but the first requires predicting the residual of
the \emph{better} predictor. A stronger predictor has already absorbed the
systematic part of that residual, so what remains is closer to noise, and the
expected ranking is increasingly governed by the response penalty alone. This
is a boundary condition of utility routing rather than a defect of a particular
router: as fixed predictors improve, context selection has less ex-ante
identifiable value, even while its oracle value persists. Practically, the
diagnostic to run before deploying a router is the one used here: compare the
oracle ladder against the best response-free ranking and check whether the
cached view correlates with the true marginal at all. Both comparisons are
computed from the deployment's own protocol and cost one sweep each. Fitting a
predictability probe is deliberately \emph{not} on this list: as
\cref{sec:strong-backbone} shows, even at $2{,}000$ training episodes the probe
cannot separate a weak expert from a strong one, so it can rank geometries but
not cells.

\subsection{Where sequential routing pays off}

\Cref{tab:boundary} collects the settings evaluated in this paper. Two
conditions must hold together. The predictor must leave a systematic residual
for contexts to correct: the ridge expert does, and the learned router converts
a double-digit share of the standalone oracle into gain with intervals well
above zero. The candidate set must make a candidate's value depend on what is
already selected: constructed redundant traffic pools do, whereas
demonstration pools without redundancy have a flat state-conditioned ceiling
($H_{\mathrm{state}}\le 5\times10^{-4}$). When the predictor is strong and
well calibrated, the first condition weakens even though the oracle gap
remains, and no evaluated method --- utility-based or policy-gradient ---
recovers it. When the pools are not redundant, the second condition fails and
nearest-neighbour relevance is within a fraction of a percent of the oracle.
The stopping rule is a separate design axis: forcing the budget to be spent
turns a 3-of-5 result into 5-of-5 in the second family without changing the
ranking.

\begin{table}[t]
 \centering
 \footnotesize
 \caption{Where the framework pays off. $\Delta$Static is R-MUR at $q=4$
 against Standalone Utility with a target-level interval.}
 \label{tab:boundary}
 \begin{tabular}{lcc}
  \toprule
  Setting & Headroom & $\Delta$Static\\
  \midrule
  Traffic, ridge (32 targets) & .09--.11 & $+.0122\,[+.010,+.015]$\\
  Traffic, ridge (PEMS-BAY) & .11 & $+.0288\,[+.021,+.040]$\\
  Traffic, nonlinear & .15 & $+.0028\,[-.001,+.008]$\\
  Traffic, nonlinear, unadapted & .16 & $+.0122\,[+.006,+.020]$\\
  PEMS-BAY, nonlinear & .15 & $-.0008\,[-.004,+.002]$\\
  Demos, in-context & $\le 5\cdot10^{-4}$ & 3/5 benchmarks\\
  Demos, forced budget & $\le 5\cdot10^{-4}$ & 5/5 benchmarks\\
  Demos, redundant pools & $\le 9\cdot10^{-4}$ & 1/5; 5/5 forced\\
  \bottomrule
 \end{tabular}
\end{table}

Two further axes decide whether routing is worth attempting, and both are
counter-intuitive. The first is adaptation. Fine-tuning the fixed predictor on
the deployment is normally free accuracy, and it is exactly the operation that
destroys routing value: in a controlled sweep of adaptation epochs, two epochs
improve anchor error by $7.1\%$ while removing about five sixths of the share
of the oracle the router could realise, and the structural opportunity is
unchanged throughout ($+.0831$ to $+.0858$). Twenty-eight further epochs buy a
further $1.2\%$ of accuracy and no identifiable value at all. The practical
reading is that routing and heavy test-time adaptation are substitutes: the
residual a router exploits is the same residual adaptation absorbs, so a
practitioner should decide which of the two is to capture it rather than
expecting both. The second axis is the size of the candidate library. Enlarging
the pool from sixteen to one hundred and twenty-eight candidates raises the
standalone-oracle opportunity by forty percent and drives the router to the
level of random selection, because the value that appears is value nothing can
identify in advance. Retrieving from a bigger library therefore makes this
problem harder, not easier --- the opposite of the usual scaling intuition ---
and the useful move at large pool sizes is to make the pool more redundant or
the predictor less adapted, not larger.

\subsection{When a simple baseline is hard to beat}

The second task family shows the same boundary from the other side. Because
the demonstration pools are not redundant, the sequential and standalone
oracles coincide and a nearest-neighbour relevance rule reaches
$99.4$--$99.9\%$ of the oracle gain. In that regime the ranking problem is
essentially solved by similarity, and no learned router, including R-MUR,
improves on it uniformly. The framework is therefore best applied where
candidates overlap in what they contribute: redundancy is what makes the value
of a candidate depend on what has already been selected, and that dependence is
what a sequential router can exploit.
